\section{Justificativa}

Os algoritmos genéticos (AGs) apresentam-se como uma abordagem poderosa e flexível para a exploração de grandes espaços de solução. Inspirados nos processos de seleção natural, esses algoritmos são capazes de lidar com problemas onde métodos exatos se tornam inviáveis devido ao crescimento exponencial da complexidade. O uso de AGs no problema das \textit{n}-rainhas permite não apenas a identificação de soluções válidas, mas também a experimentação com diferentes parâmetros evolutivos, proporcionando uma rica base para análise e compreensão do comportamento dessas técnicas.

A escolha do problema das \textit{n}-rainhas como objeto central desta pesquisa se fundamenta em sua relevância histórica, pedagógica e prática na área de computação. Trata-se de um problema clássico que, apesar de sua formulação simples, apresenta um grau de complexidade combinatória crescente à medida que o número de rainhas aumenta. Uma de suas principais vantagens para estudos experimentais é a facilidade de replicação: para um valor fixo de \textit{n}, o problema permanece invariável — por exemplo, o cenário com 8 rainhas é sempre o mesmo — ao contrário de outros problemas combinatórios com \textit{n} elementos, em que diferentes instâncias podem alterar significativamente o espaço de busca. Essa característica torna o problema das \textit{n}-rainhas particularmente adequado como benchmark, permitindo análises controladas, reprodutíveis e comparáveis entre diferentes abordagens algorítmicas. Por essas razões, o problema é amplamente utilizado no estudo e aplicação de algoritmos de busca, heurísticas e métodos de otimização.

Este trabalho propõe investigar de forma sistemática o impacto de diferentes estratégias de controle populacional. A escolha por essa abordagem se justifica pela escassez de estudos que realizem uma comparação direta entre essas variações, especialmente em um cenário bem definido como o das \textit{n}-rainhas. Ao manter constante o problema e a estrutura geral do algoritmo, torna-se possível isolar e compreender melhor o efeito de cada escolha de configuração sobre o desempenho do sistema.

A implementação prática em \texttt{Python}, combinada à execução de experimentos com diferentes valores de \textit{n}, favorece a replicabilidade em contextos acadêmicos. Como linguagem de alto nível, legível e amplamente adotada nas áreas de ciência de dados e inteligência artificial, o uso de Python torna os experimentos mais acessíveis, incentivando a reprodutibilidade e o reaproveitamento do código em pesquisas futuras.

Por fim, esta pesquisa é relevante por promover a compreensão prática dos algoritmos genéticos, incentivando a análise crítica de suas configurações e demonstrando como escolhas aparentemente simples — como a forma de controlar os indivíduos — podem afetar significativamente os resultados. Ao aprofundar-se nessa análise, o trabalho contribui para a consolidação de boas práticas na aplicação de algoritmos evolutivos e amplia a compreensão sobre o potencial dessas técnicas na resolução de problemas complexos.

